\documentclass[11pt, a4paper]{article}
\usepackage[utf8]{inputenc}
\usepackage[english]{babel}
\usepackage{amsmath}
\usepackage{amssymb}
\usepackage{csquotes}
\usepackage[style=authoryear, backend=biber]{biblatex}
\usepackage{setspace}
\usepackage{parskip}
\usepackage{enumitem}
\usepackage{xcolor}
\usepackage{tikz}
\usetikzlibrary{positioning}

\title{Autonomous Interdependence and the Evaluation of Human Relationships}
\author{}
\date{}

\begin{document}

\maketitle

\section*{Introduction}
Human relationships can be approached as systems of interaction between autonomous but mutually dependent organisms.

The two terms are important. \textbf{Autonomy} implies that individuals possess distinct internal states, preferences, objectives, and capacities for action. \textbf{Interdependence} implies that important components of their functioning nevertheless depend on the behaviour of others. Human beings are therefore neither fully independent agents nor components of a single collective organism. They remain separate while continuously affecting, supporting, constraining, and modifying one another.

This produces a general problem of coordination.

\section*{The Problem of Coordination}
Each individual acts from a perspective that is only partially accessible to others. Internal states—such as preference, discomfort, affection, fear, desire, fatigue, and expectation—cannot be directly observed. They must be communicated, inferred, and continuously updated. At the same time, decisions taken by one individual alter the conditions experienced by another.

Human relationships can consequently be described as \emph{repeated processes of mutual adjustment under conditions of incomplete information}.

\section*{Form vs. Quality}
From this perspective, a distinction can be made between the \textbf{form of a relationship} and the \textbf{quality of the interactions occurring within it}.

\begin{itemize}[leftmargin=*]
    \item \textbf{Relationship form} concerns its socially recognisable configuration: marriage, monogamy, cohabitation, friendship, kinship, sexual partnership, or another interpersonal arrangement.
    \item \textbf{Interaction quality} concerns what actually occurs between the individuals participating in that configuration: whether information can be exchanged, preferences expressed, boundaries maintained, burdens distributed, disagreements negotiated, mistakes corrected, and support reciprocated.
\end{itemize}

There is no necessary equivalence between these two dimensions.

A socially recognised or culturally preferred relationship structure does not by itself guarantee functional interaction. Conversely, valuable relational functions may occur outside the conventional structure to which they are commonly assigned. Emotional support, practical assistance, intellectual companionship, caregiving, physical intimacy, and sexual interaction are distinct relational functions. There is no structural requirement that all of them must be supplied by the same individual.

\section*{A General Proposition}
This leads to a general proposition:

\begin{quote}
    \textbf{Autonomous interdependence does not prescribe the shape of a good relationship; it constrains the quality of the interaction.}
\end{quote}

Under this proposition, the principal object of evaluation is not whether a relationship conforms to a predetermined interpersonal architecture. It is whether the interactions occurring within that architecture remain compatible with both autonomy and interdependence.

This immediately excludes two limiting cases.

\begin{itemize}[leftmargin=*]
    \item At one extreme lies \textbf{unilateral optimisation}. If one participant systematically maximises personal benefit while transferring costs to another, the interaction becomes extractive. Interdependence remains, but autonomy and reciprocity are progressively reduced.
    \item At the opposite extreme lies \textbf{unilateral self-sacrifice}. Here, one participant continuously subordinates personal interests, preferences, or boundaries to those of another. Although such behaviour may superficially resemble care, it can similarly produce an asymmetric system in which one participant's functioning becomes dependent on the progressive reduction of the other's autonomy.
\end{itemize}

Neither configuration provides a stable model of autonomous interdependence.

\section*{The Relational Problem}
The relevant relational problem therefore exists between these extremes. Participants must continuously negotiate how much to give, receive, request, refuse, provide, and depend upon one another.

There is no universal equilibrium.

The appropriate distribution changes over time. Illness may temporarily increase dependence. Employment may alter available time. Parenthood changes responsibilities. Stress may reduce emotional capacity. Ageing modifies physical dependence. Desire fluctuates. External circumstances change. One individual may temporarily contribute more in one domain while receiving more in another.

Consequently, a functional relationship cannot be defined solely by a fixed allocation of responsibilities.

Its more important property is \textbf{adaptability}.

\section*{Adaptability as a Key Property}
A viable relational system must be capable of detecting changes in the states of its participants and modifying its behaviour accordingly. This requires information to remain available. Needs must be expressible. Discomfort must be reportable. Preferences must be revisable. Disagreement must not automatically threaten the existence of the relationship. Participants must be able to discover that their assumptions about one another were incorrect and alter their behaviour in response.

Relationship quality is therefore partly a property of the system's capacity for correction.

This becomes particularly important because most relationships consist not of exceptional moments but of repeated ordinary interactions.

Food is prepared. Time is allocated. Questions are answered. Attention is given or withheld. Responsibilities are remembered or forgotten. Emotional signals are noticed or missed. Plans are negotiated. Preferences conflict. Minor injuries occur and are repaired. Information is shared. Decisions are revised.

Individually, many of these events are trivial. Collectively, they constitute the relationship.

The experiential quality of a relationship therefore emerges substantially from the accumulated structure of everyday interaction. A relationship may contain occasional exceptional experiences while being chronically poor in its ordinary exchanges. Conversely, a relationship consisting largely of mundane activities may nevertheless generate substantial wellbeing when those interactions reliably preserve safety, recognition, reciprocity, and mutual responsiveness.

From this viewpoint, relational quality is not primarily a property of intensity. It is a property of repeated coordination.

\section*{Relational Plurality}
This also changes the significance of relational plurality.

If relationship quality is evaluated through interaction rather than exclusively through form, there is no theoretical requirement that all important forms of dependence must be concentrated in a single interpersonal bond. Human beings already function within networks of partially overlapping dependencies. Different individuals provide different forms of information, assistance, companionship, affection, care, and practical capacity.

This does not imply that individuals are interchangeable.

Historical attachment, accumulated trust, familiarity, compatibility, and shared experience create relationship-specific value. The loss of one person cannot necessarily be compensated by the addition of another. However, the functions occurring within human relationships are analytically separable from the socially defined category of the relationship itself.

The distinction matters because it shifts the normative question.

Instead of asking primarily:
\begin{quote}
    \emph{What kind of relationship should people have?}
\end{quote}

one can ask:
\begin{quote}
    \emph{What properties must interactions possess if autonomous individuals are to depend upon one another without progressively damaging either autonomy or mutual functioning?}
\end{quote}

\section*{Sexual Interaction as an Illustration}
Sexual interaction provides an unusually clear illustration of this general problem.

Sexual pleasure is experienced individually. One nervous system cannot directly experience the sensations occurring within another. Yet during interpersonal sexual activity, the experience of each participant is substantially affected by the behaviour of the other.

Sex therefore creates a particularly immediate form of interdependence under incomplete information.

Neither participant possesses direct knowledge of what the other is experiencing. Desire, pleasure, discomfort, hesitation, excitement, and boundaries must be communicated or inferred. Moreover, these states can change during the interaction itself.

A successful sexual interaction consequently requires continuous reciprocal adjustment.

Each participant simultaneously produces actions and receives information about their effects. Touch, movement, language, hesitation, proximity, and withdrawal become part of a feedback process through which the participants attempt to coordinate two private experiences.

Safety and acceptance are therefore not merely external moral requirements placed upon sexual behaviour. They are functional properties of the interaction.

If an individual cannot communicate discomfort, uncertainty, or preference, information necessary for coordination is lost. If the other participant ignores, suppresses, or fails to respond to such information, the system loses its capacity for mutual adjustment.

Sexual interaction therefore makes explicit what is less immediately visible in many other relationships: another person's internal state cannot be assumed. It must remain discoverable.

The interaction is consequently neither straightforwardly altruistic nor straightforwardly self-interested. Each individual can legitimately pursue personal pleasure, but the possibility of obtaining that pleasure depends partly upon maintaining conditions in which the other individual can also participate as an autonomous agent.

The structure is reciprocal without requiring symmetry at every moment.

\section*{Beyond Sexuality}
The same principle extends beyond sexuality.

Human relational systems function through variable mixtures of dependence, autonomy, information exchange, and adaptation. Their stability does not require that every interaction be equal, nor that all benefits be simultaneously reciprocal. It requires that asymmetries remain negotiable and that the system retain the capacity to recognise and correct forms of dependence that become persistently destructive.

The resulting relational world is necessarily untidy.

There is no reason to expect a system involving multiple autonomous agents, changing environments, private internal states, and overlapping dependencies to converge upon a single static interpersonal arrangement.

Messiness is therefore not necessarily evidence of relational failure. It is partly an expected consequence of adaptive interdependence.

What can be specified more strictly are the conditions under which such messiness remains viable.

\begin{itemize}[leftmargin=*]
    \item Information must remain sufficiently available for mutual adjustment.
    \item Boundaries must remain communicable.
    \item Dependence must not systematically eliminate autonomy.
    \item Participants must retain meaningful capacity to refuse, renegotiate, and correct.
    \item Persistent extraction must remain identifiable.
    \item Changes in circumstance must be capable of producing changes in arrangement.
\end{itemize}

These are constraints on interaction rather than prescriptions of form.

\section*{An Unusual Relational Ethic}
Autonomous interdependence therefore produces an unusual type of relational ethic. It is comparatively permissive regarding the architecture of human relationships but considerably less permissive regarding the processes occurring inside them.

Who relates to whom, in what social configuration, may admit substantial variation.

Whether the resulting interactions preserve information, autonomy, reciprocity, adaptation, and the possibility of correction is much less arbitrary.

The objective is therefore neither independence nor complete dependence.

It is the maintenance of relationships in which distinct individuals can meaningfully depend upon one another while remaining sufficiently autonomous to communicate, revise, refuse, contribute, and change.

\section*{Conclusion}
Under this formulation, the central question concerning a human relationship is not what the relationship is called.

It is what the interaction allows its participants to remain.

\end{document}